% Options for packages loaded elsewhere
\PassOptionsToPackage{unicode}{hyperref}
\PassOptionsToPackage{hyphens}{url}
\documentclass[
  11pt,
  a4paper,
]{article}
\usepackage{xcolor}
\usepackage[margin=18mm]{geometry}
\usepackage{amsmath,amssymb}
\setcounter{secnumdepth}{-\maxdimen} % remove section numbering
\usepackage{iftex}
\ifPDFTeX
  \usepackage[T1]{fontenc}
  \usepackage[utf8]{inputenc}
  \usepackage{textcomp} % provide euro and other symbols
\else % if luatex or xetex
  \usepackage{unicode-math} % this also loads fontspec
  \defaultfontfeatures{Scale=MatchLowercase}
  \defaultfontfeatures[\rmfamily]{Ligatures=TeX,Scale=1}
\fi
\ifPDFTeX\else
  % xetex/luatex font selection
    \setmainfont[]{Noto Serif}
    \setsansfont[]{Noto Sans}
    \setmonofont[]{Noto Sans Mono}
  \setmathfont[]{Noto Sans Math}
\fi
% Use upquote if available, for straight quotes in verbatim environments
\IfFileExists{upquote.sty}{\usepackage{upquote}}{}
\IfFileExists{microtype.sty}{% use microtype if available
  \usepackage[]{microtype}
  \UseMicrotypeSet[protrusion]{basicmath} % disable protrusion for tt fonts
}{}
\makeatletter
\@ifundefined{KOMAClassName}{% if non-KOMA class
  \IfFileExists{parskip.sty}{%
    \usepackage{parskip}
  }{% else
    \setlength{\parindent}{0pt}
    \setlength{\parskip}{6pt plus 2pt minus 1pt}}
}{% if KOMA class
  \KOMAoptions{parskip=half}}
\makeatother
\ifLuaTeX
\usepackage[bidi=basic]{babel}
\else
\usepackage[bidi=default]{babel}
\fi
\babelprovide[main,import]{russian}
\ifPDFTeX
\else
\babelfont{rm}[]{Noto Serif}
\fi
% get rid of language-specific shorthands (see #6817):
\let\LanguageShortHands\languageshorthands
\def\languageshorthands#1{}
\setlength{\emergencystretch}{3em} % prevent overfull lines
\providecommand{\tightlist}{%
  \setlength{\itemsep}{0pt}\setlength{\parskip}{0pt}}
\usepackage{bookmark}
\IfFileExists{xurl.sty}{\usepackage{xurl}}{} % add URL line breaks if available
\urlstyle{same}
\hypersetup{
  pdftitle={Билеты по дискретным случайным процессам},
  pdflang={ru-RU},
  hidelinks,
  pdfcreator={LaTeX via pandoc}}

\title{Билеты по дискретным случайным процессам}
\author{}
\date{}

\begin{document}
\maketitle

\section{Билет 02. Счетная аддитивность и вероятностные пространства.
Мера, порожденная функцией
распределения}\label{ux431ux438ux43bux435ux442-02.-ux441ux447ux435ux442ux43dux430ux44f-ux430ux434ux434ux438ux442ux438ux432ux43dux43eux441ux442ux44c-ux438-ux432ux435ux440ux43eux44fux442ux43dux43eux441ux442ux43dux44bux435-ux43fux440ux43eux441ux442ux440ux430ux43dux441ux442ux432ux430.-ux43cux435ux440ux430-ux43fux43eux440ux43eux436ux434ux435ux43dux43dux430ux44f-ux444ux443ux43dux43aux446ux438ux435ux439-ux440ux430ux441ux43fux440ux435ux434ux435ux43bux435ux43dux438ux44f}

Источники: Ширяев 1, гл. II; лекции ДСП.

\subsection{1. Короткий
ответ}\label{ux43aux43eux440ux43eux442ux43aux438ux439-ux43eux442ux432ux435ux442}

Вероятностное пространство в аксиоматике Колмогорова - это тройка

\[
(\Omega,\mathcal F,\mathbb P),
\]

где \(\Omega\) - пространство элементарных исходов, \(\mathcal F\) -
\(\sigma\)-алгебра событий, а \(\mathbb P\) - счетно-аддитивная мера на
\(\mathcal F\) с полной массой \(1\).

Главная аксиома:

\[
A_i\cap A_j=\varnothing,\ i\ne j
\quad\Longrightarrow\quad
\mathbb P\left(\bigsqcup_{n=1}^{\infty} A_n\right)
=
\sum_{n=1}^{\infty}\mathbb P(A_n).
\]

Из нее следуют \(\mathbb P(\varnothing)=0\), конечная аддитивность,
формула дополнения, монотонность, счетная субаддитивность и
непрерывность вероятности снизу и сверху.

Если \(F:\mathbb R\to[0,1]\) не убывает, непрерывна справа и

\[
\lim_{x\to-\infty}F(x)=0,\qquad
\lim_{x\to+\infty}F(x)=1,
\]

то \(F\) порождает единственную вероятностную меру \(\mu_F\) на
борелевской \(\sigma\)-алгебре \(\mathcal B(\mathbb R)\) такую, что

\[
\mu_F((a,b])=F(b)-F(a),\qquad a<b.
\]

Это мера Лебега-Стилтьеса, соответствующая функции распределения \(F\).
В этом билете существование такой меры обычно формулируется без
доказательства; доказательство опирается на теорему продолжения меры.

\subsection{2. Полный
ответ}\label{ux43fux43eux43bux43dux44bux439-ux43eux442ux432ux435ux442}

\subsubsection{2.1. Аксиоматический подход
Колмогорова}\label{ux430ux43aux441ux438ux43eux43cux430ux442ux438ux447ux435ux441ux43aux438ux439-ux43fux43eux434ux445ux43eux434-ux43aux43eux43bux43cux43eux433ux43eux440ux43eux432ux430}

В первом билете выбиралась система множеств, на которой можно говорить о
событиях. Во втором билете на этой системе задается вероятность. Идея
Колмогорова состоит в том, что вероятность является обычной мерой, но
нормированной условием полной массы:

\[
\mathbb P(\Omega)=1.
\]

Пусть \(\Omega\) - множество всех элементарных исходов эксперимента.
Например, для одного броска монеты \(\Omega=\{H,T\}\), для бесконечной
последовательности бросков \(\Omega=\{H,T\}^{\mathbb N}\), для
вещественной случайной величины естественно брать \(\Omega=\mathbb R\).
Не каждое подмножество \(\Omega\) автоматически объявляется событием.
Выбирается \(\sigma\)-алгебра \(\mathcal F\subset 2^\Omega\), то есть
класс подмножеств, замкнутый относительно дополнений и счетных
объединений. Элементы \(\mathcal F\) называются событиями.

\subsubsection{2.2. Определение и аксиомы
вероятности}\label{ux43eux43fux440ux435ux434ux435ux43bux435ux43dux438ux435-ux438-ux430ux43aux441ux438ux43eux43cux44b-ux432ux435ux440ux43eux44fux442ux43dux43eux441ux442ux438}

Вероятность - это функция

\[
\mathbb P:\mathcal F\to[0,1],
\]

которая удовлетворяет трем требованиям:

\begin{enumerate}
\def\labelenumi{\arabic{enumi}.}
\tightlist
\item
  неотрицательность: \(\mathbb P(A)\ge 0\) для всех \(A\in\mathcal F\);
\item
  нормировка: \(\mathbb P(\Omega)=1\);
\item
  счетная аддитивность: если \(A_1,A_2,\ldots\in\mathcal F\) попарно не
  пересекаются, то
\end{enumerate}

\[
\mathbb P\left(\bigcup_{n=1}^{\infty}A_n\right)
=
\sum_{n=1}^{\infty}\mathbb P(A_n).
\]

Так как множества \(A_n\) попарно не пересекаются, объединение часто
обозначают дизъюнктным объединением:

\[
\bigcup_{n=1}^{\infty}A_n=\bigsqcup_{n=1}^{\infty}A_n.
\]

\subsubsection{2.3. Счетная и конечная
аддитивность}\label{ux441ux447ux435ux442ux43dux430ux44f-ux438-ux43aux43eux43dux435ux447ux43dux430ux44f-ux430ux434ux434ux438ux442ux438ux432ux43dux43eux441ux442ux44c}

Счетная аддитивность является принципиальным усилением конечной
аддитивности. Конечная аддитивность говорит только о конечных
разбиениях:

\[
\mathbb P\left(\bigsqcup_{k=1}^m A_k\right)
=
\sum_{k=1}^m\mathbb P(A_k).
\]

Для вероятностной теории этого недостаточно, потому что основные события
часто являются счетными объединениями или пересечениями. Например,
событие ``хотя бы однажды произошел успех'' имеет вид

\[
\bigcup_{n=1}^{\infty}A_n,
\]

событие ``успехи происходят бесконечно часто'' имеет вид

\[
\limsup_{n\to\infty}A_n
=
\bigcap_{m=1}^{\infty}\bigcup_{n=m}^{\infty}A_n,
\]

а предельные теоремы работают именно с пределами последовательностей
событий и случайных величин. Поэтому в аксиоматике Колмогорова требуется
не просто конечная, а счетная аддитивность.

\subsubsection{2.4. Свойства
вероятности}\label{ux441ux432ux43eux439ux441ux442ux432ux430-ux432ux435ux440ux43eux44fux442ux43dux43eux441ux442ux438}

Из аксиом вероятности выводится набор стандартных свойств:
\(\mathbb P(\varnothing)=0\), формула дополнения, монотонность, формула
сложения для объединения, субаддитивность, непрерывность снизу и
непрерывность сверху. Эти свойства не являются дополнительными
аксиомами. Их нужно уметь быстро доказывать из счетной аддитивности.

\subsubsection{2.5. Вероятностная мера и функция
распределения}\label{ux432ux435ux440ux43eux44fux442ux43dux43eux441ux442ux43dux430ux44f-ux43cux435ux440ux430-ux438-ux444ux443ux43dux43aux446ux438ux44f-ux440ux430ux441ux43fux440ux435ux434ux435ux43bux435ux43dux438ux44f}

Вторая часть билета - пример вероятностной меры, заданной не через явное
перечисление исходов, а через функцию распределения. Пусть
\(F:\mathbb R\to[0,1]\) обладает свойствами:

\begin{enumerate}
\def\labelenumi{\arabic{enumi}.}
\tightlist
\item
  \(F\) не убывает;
\item
  \(F\) непрерывна справа;
\item
  \(\lim_{x\to-\infty}F(x)=0\);
\item
  \(\lim_{x\to+\infty}F(x)=1\).
\end{enumerate}

Тогда существует единственная вероятностная мера \(\mu_F\) на
\(\mathcal B(\mathbb R)\), для которой

\[
\mu_F((-\infty,x])=F(x)
\]

и, следовательно,

\[
\mu_F((a,b])=F(b)-F(a).
\]

\subsubsection{2.6. Мера Лебега-Стилтьеса и типы
распределений}\label{ux43cux435ux440ux430-ux43bux435ux431ux435ux433ux430-ux441ux442ux438ux43bux442ux44cux435ux441ux430-ux438-ux442ux438ux43fux44b-ux440ux430ux441ux43fux440ux435ux434ux435ux43bux435ux43dux438ux439}

Эта мера называется мерой Лебега-Стилтьеса, порожденной \(F\).
Интуитивно \(F(x)\) задает массу луча \((-\infty,x]\), а разность
\(F(b)-F(a)\) задает массу полуинтервала \((a,b]\). Если \(F\) имеет
скачок в точке \(a\), то в этой точке у меры есть атом:

\[
\mu_F(\{a\})=F(a)-F(a-),
\]

где

\[
F(a-)=\lim_{x\uparrow a}F(x).
\]

Если \(F\) непрерывна, атомов нет. Если \(F\) абсолютно непрерывна и
имеет плотность \(f\), то

\[
F(x)=\int_{-\infty}^{x} f(t)\,dt,\qquad
\mu_F((a,b])=\int_a^b f(t)\,dt.
\]

Таким образом, дискретные, непрерывные и смешанные распределения входят
в одну общую схему: каждое распределение на прямой задается своей
функцией распределения.

\subsection{3. Основные
определения}\label{ux43eux441ux43dux43eux432ux43dux44bux435-ux43eux43fux440ux435ux434ux435ux43bux435ux43dux438ux44f}

\subsubsection{Измеримое
пространство}\label{ux438ux437ux43cux435ux440ux438ux43cux43eux435-ux43fux440ux43eux441ux442ux440ux430ux43dux441ux442ux432ux43e}

Пара \((\Omega,\mathcal F)\), где \(\Omega\) - множество, а
\(\mathcal F\) - \(\sigma\)-алгебра подмножеств \(\Omega\).

\subsubsection{Элементарный
исход}\label{ux44dux43bux435ux43cux435ux43dux442ux430ux440ux43dux44bux439-ux438ux441ux445ux43eux434}

Элемент \(\omega\in\Omega\). Он описывает один возможный исход
эксперимента.

\subsubsection{Событие}\label{ux441ux43eux431ux44bux442ux438ux435}

Множество \(A\in\mathcal F\). Если в эксперименте реализовался исход
\(\omega\in A\), говорят, что событие \(A\) произошло.

\subsubsection{Вероятностная
мера}\label{ux432ux435ux440ux43eux44fux442ux43dux43eux441ux442ux43dux430ux44f-ux43cux435ux440ux430}

Функция \(\mathbb P:\mathcal F\to[0,1]\), удовлетворяющая условиям:

\[
\mathbb P(\Omega)=1,
\]

и для любых попарно непересекающихся \(A_1,A_2,\ldots\in\mathcal F\):

\[
\mathbb P\left(\bigsqcup_{n=1}^{\infty}A_n\right)
=
\sum_{n=1}^{\infty}\mathbb P(A_n).
\]

Неотрицательность обычно включается в определение меры. Так как значения
берутся в \([0,1]\), она уже зафиксирована.

\subsubsection{Вероятностное пространство, или абстрактная
колмогоровская
тройка}\label{ux432ux435ux440ux43eux44fux442ux43dux43eux441ux442ux43dux43eux435-ux43fux440ux43eux441ux442ux440ux430ux43dux441ux442ux432ux43e-ux438ux43bux438-ux430ux431ux441ux442ux440ux430ux43aux442ux43dux430ux44f-ux43aux43eux43bux43cux43eux433ux43eux440ux43eux432ux441ux43aux430ux44f-ux442ux440ux43eux439ux43aux430}

Тройка

\[
(\Omega,\mathcal F,\mathbb P),
\]

где \((\Omega,\mathcal F)\) - измеримое пространство, а \(\mathbb P\) -
вероятностная мера на \(\mathcal F\).

\subsubsection{Попарно несовместные
события}\label{ux43fux43eux43fux430ux440ux43dux43e-ux43dux435ux441ux43eux432ux43cux435ux441ux442ux43dux44bux435-ux441ux43eux431ux44bux442ux438ux44f}

Последовательность \((A_n)_{n\ge1}\) называется попарно несовместной,
если

\[
A_i\cap A_j=\varnothing,\qquad i\ne j.
\]

\subsubsection{Конечная
аддитивность}\label{ux43aux43eux43dux435ux447ux43dux430ux44f-ux430ux434ux434ux438ux442ux438ux432ux43dux43eux441ux442ux44c}

Для попарно несовместных \(A_1,\ldots,A_m\in\mathcal F\):

\[
\mathbb P\left(\bigsqcup_{k=1}^{m}A_k\right)
=
\sum_{k=1}^{m}\mathbb P(A_k).
\]

\subsubsection{Счетная
аддитивность}\label{ux441ux447ux435ux442ux43dux430ux44f-ux430ux434ux434ux438ux442ux438ux432ux43dux43eux441ux442ux44c}

Для попарно несовместных \(A_1,A_2,\ldots\in\mathcal F\):

\[
\mathbb P\left(\bigsqcup_{k=1}^{\infty}A_k\right)
=
\sum_{k=1}^{\infty}\mathbb P(A_k).
\]

\subsubsection{\texorpdfstring{Борелевская \(\sigma\)-алгебра на
\(\mathbb R\)}{Борелевская \textbackslash sigma-алгебра на \textbackslash mathbb R}}\label{ux431ux43eux440ux435ux43bux435ux432ux441ux43aux430ux44f-sigma-ux430ux43bux433ux435ux431ux440ux430-ux43dux430-mathbb-r}

Минимальная \(\sigma\)-алгебра, содержащая все открытые множества
прямой. Обозначается \(\mathcal B(\mathbb R)\).

Эквивалентно, \(\mathcal B(\mathbb R)\) порождена интервалами вида
\((-\infty,x]\):

\[
\mathcal B(\mathbb R)=\sigma\{(-\infty,x]:x\in\mathbb R\}.
\]

\subsubsection{\texorpdfstring{Функция распределения меры \(\mu\) на
\(\mathbb R\)}{Функция распределения меры \textbackslash mu на \textbackslash mathbb R}}\label{ux444ux443ux43dux43aux446ux438ux44f-ux440ux430ux441ux43fux440ux435ux434ux435ux43bux435ux43dux438ux44f-ux43cux435ux440ux44b-mu-ux43dux430-mathbb-r}

\[
F_\mu(x)=\mu((-\infty,x]).
\]

\subsubsection{\texorpdfstring{Функция распределения случайной величины
\(X\)}{Функция распределения случайной величины X}}\label{ux444ux443ux43dux43aux446ux438ux44f-ux440ux430ux441ux43fux440ux435ux434ux435ux43bux435ux43dux438ux44f-ux441ux43bux443ux447ux430ux439ux43dux43eux439-ux432ux435ux43bux438ux447ux438ux43dux44b-x}

\[
F_X(x)=\mathbb P(X\le x)
=
\mathbb P(\{\omega:X(\omega)\le x\}).
\]

\subsubsection{Абстрактная функция
распределения}\label{ux430ux431ux441ux442ux440ux430ux43aux442ux43dux430ux44f-ux444ux443ux43dux43aux446ux438ux44f-ux440ux430ux441ux43fux440ux435ux434ux435ux43bux435ux43dux438ux44f}

Функция \(F:\mathbb R\to[0,1]\), которая не убывает, непрерывна справа и
имеет пределы \(0\) на \(-\infty\) и \(1\) на \(+\infty\).

\subsubsection{Мера, порожденная функцией
распределения}\label{ux43cux435ux440ux430-ux43fux43eux440ux43eux436ux434ux435ux43dux43dux430ux44f-ux444ux443ux43dux43aux446ux438ux435ux439-ux440ux430ux441ux43fux440ux435ux434ux435ux43bux435ux43dux438ux44f}

Единственная вероятностная мера \(\mu_F\) на \(\mathcal B(\mathbb R)\),
для которой

\[
\mu_F((-\infty,x])=F(x)
\]

для всех \(x\in\mathbb R\). Тогда для \(a<b\):

\[
\mu_F((a,b])=F(b)-F(a).
\]

\subsection{4. Основные теоремы и
свойства}\label{ux43eux441ux43dux43eux432ux43dux44bux435-ux442ux435ux43eux440ux435ux43cux44b-ux438-ux441ux432ux43eux439ux441ux442ux432ux430}

\subsubsection{Свойство 1. Вероятность пустого множества равна
нулю}\label{ux441ux432ux43eux439ux441ux442ux432ux43e-1.-ux432ux435ux440ux43eux44fux442ux43dux43eux441ux442ux44c-ux43fux443ux441ux442ux43eux433ux43e-ux43cux43dux43eux436ux435ux441ux442ux432ux430-ux440ux430ux432ux43dux430-ux43dux443ux43bux44e}

\[
\mathbb P(\varnothing)=0.
\]

\subsubsection{Доказательство}\label{ux434ux43eux43aux430ux437ux430ux442ux435ux43bux44cux441ux442ux432ux43e}

Используем дизъюнктное разложение

\[
\Omega=\Omega\sqcup\varnothing\sqcup\varnothing\sqcup\cdots.
\]

По счетной аддитивности

\[
1=\mathbb P(\Omega)
=
\mathbb P(\Omega)+\sum_{n=2}^{\infty}\mathbb P(\varnothing)
=
1+\sum_{n=2}^{\infty}\mathbb P(\varnothing).
\]

Следовательно,

\[
\sum_{n=2}^{\infty}\mathbb P(\varnothing)=0.
\]

Слагаемые неотрицательны, значит \(\mathbb P(\varnothing)=0\).
\(\square\)

\subsubsection{Свойство 2. Из счетной аддитивности следует конечная
аддитивность}\label{ux441ux432ux43eux439ux441ux442ux432ux43e-2.-ux438ux437-ux441ux447ux435ux442ux43dux43eux439-ux430ux434ux434ux438ux442ux438ux432ux43dux43eux441ux442ux438-ux441ux43bux435ux434ux443ux435ux442-ux43aux43eux43dux435ux447ux43dux430ux44f-ux430ux434ux434ux438ux442ux438ux432ux43dux43eux441ux442ux44c}

Если \(A_1,\ldots,A_m\) попарно не пересекаются, то

\[
\mathbb P\left(\bigsqcup_{k=1}^{m}A_k\right)
=
\sum_{k=1}^{m}\mathbb P(A_k).
\]

\subsubsection{Доказательство}\label{ux434ux43eux43aux430ux437ux430ux442ux435ux43bux44cux441ux442ux432ux43e-1}

Возьмем счетную последовательность

\[
A_1,\ldots,A_m,\varnothing,\varnothing,\ldots
\]

и применим счетную аддитивность:

\[
\mathbb P\left(\bigsqcup_{k=1}^{m}A_k\right)
=
\sum_{k=1}^{m}\mathbb P(A_k)+\sum_{k=m+1}^{\infty}\mathbb P(\varnothing).
\]

По свойству 1 \(\mathbb P(\varnothing)=0\), поэтому

\[
\mathbb P\left(\bigsqcup_{k=1}^{m}A_k\right)
=
\sum_{k=1}^{m}\mathbb P(A_k).
\]

\(\square\)

\subsubsection{Свойство 3. Формула
дополнения}\label{ux441ux432ux43eux439ux441ux442ux432ux43e-3.-ux444ux43eux440ux43cux443ux43bux430-ux434ux43eux43fux43eux43bux43dux435ux43dux438ux44f}

Для любого \(A\in\mathcal F\):

\[
\mathbb P(\overline{A})=1-\mathbb P(A).
\]

\subsubsection{Доказательство}\label{ux434ux43eux43aux430ux437ux430ux442ux435ux43bux44cux441ux442ux432ux43e-2}

Множества \(A\) и \(\overline{A}\) несовместны, причем

\[
A\sqcup \overline{A}=\Omega.
\]

По конечной аддитивности:

\[
1=\mathbb P(\Omega)=\mathbb P(A)+\mathbb P(\overline{A}).
\]

Отсюда

\[
\mathbb P(\overline{A})=1-\mathbb P(A).
\]

\(\square\)

\subsubsection{Свойство 4. Монотонность и формула
разности}\label{ux441ux432ux43eux439ux441ux442ux432ux43e-4.-ux43cux43eux43dux43eux442ux43eux43dux43dux43eux441ux442ux44c-ux438-ux444ux43eux440ux43cux443ux43bux430-ux440ux430ux437ux43dux43eux441ux442ux438}

Если \(A\subset B\), то

\[
\mathbb P(B\setminus A)=\mathbb P(B)-\mathbb P(A),
\qquad
\mathbb P(A)\le \mathbb P(B).
\]

\subsubsection{Доказательство}\label{ux434ux43eux43aux430ux437ux430ux442ux435ux43bux44cux441ux442ux432ux43e-3}

При \(A\subset B\) имеем дизъюнктное разложение

\[
B=A\sqcup(B\setminus A).
\]

Значит

\[
\mathbb P(B)=\mathbb P(A)+\mathbb P(B\setminus A).
\]

Отсюда следует формула разности. Так как
\(\mathbb P(B\setminus A)\ge0\), получаем монотонность:

\[
\mathbb P(A)\le\mathbb P(B).
\]

\(\square\)

\subsubsection{Свойство 5. Формула сложения для двух
событий}\label{ux441ux432ux43eux439ux441ux442ux432ux43e-5.-ux444ux43eux440ux43cux443ux43bux430-ux441ux43bux43eux436ux435ux43dux438ux44f-ux434ux43bux44f-ux434ux432ux443ux445-ux441ux43eux431ux44bux442ux438ux439}

Для любых \(A,B\in\mathcal F\):

\[
\mathbb P(A\cup B)=\mathbb P(A)+\mathbb P(B)-\mathbb P(A\cap B).
\]

\subsubsection{Доказательство}\label{ux434ux43eux43aux430ux437ux430ux442ux435ux43bux44cux441ux442ux432ux43e-4}

Разложим \(A\cup B\) на попарно непересекающиеся части:

\[
A\cup B=(A\setminus B)\sqcup(B\setminus A)\sqcup(A\cap B).
\]

Также

\[
A=(A\setminus B)\sqcup(A\cap B),
\qquad
B=(B\setminus A)\sqcup(A\cap B).
\]

По конечной аддитивности:

\[
\mathbb P(A\cup B)
=
\mathbb P(A\setminus B)+\mathbb P(B\setminus A)+\mathbb P(A\cap B),
\]

а

\[
\mathbb P(A)+\mathbb P(B)
=
\mathbb P(A\setminus B)+\mathbb P(B\setminus A)+2\mathbb P(A\cap B).
\]

Вычитая один раз \(\mathbb P(A\cap B)\), получаем нужную формулу.
\(\square\)

\subsubsection{Свойство 6. Счетная
субаддитивность}\label{ux441ux432ux43eux439ux441ux442ux432ux43e-6.-ux441ux447ux435ux442ux43dux430ux44f-ux441ux443ux431ux430ux434ux434ux438ux442ux438ux432ux43dux43eux441ux442ux44c}

Для любых событий \(A_1,A_2,\ldots\in\mathcal F\):

\[
\mathbb P\left(\bigcup_{n=1}^{\infty}A_n\right)
\le
\sum_{n=1}^{\infty}\mathbb P(A_n).
\]

\subsubsection{Доказательство}\label{ux434ux43eux43aux430ux437ux430ux442ux435ux43bux44cux441ux442ux432ux43e-5}

Преобразуем произвольную последовательность событий в попарно
непересекающуюся последовательность:

\[
B_1=A_1,\qquad
B_n=A_n\setminus\bigcup_{k=1}^{n-1}A_k,\quad n\ge2.
\]

Тогда \(B_n\subset A_n\), множества \(B_n\) попарно не пересекаются и

\[
\bigcup_{n=1}^{\infty}A_n
=
\bigsqcup_{n=1}^{\infty}B_n.
\]

По счетной аддитивности и монотонности:

\[
\mathbb P\left(\bigcup_{n=1}^{\infty}A_n\right)
=
\sum_{n=1}^{\infty}\mathbb P(B_n)
\le
\sum_{n=1}^{\infty}\mathbb P(A_n).
\]

\(\square\)

\subsubsection{Свойство 7. Непрерывность вероятности
снизу}\label{ux441ux432ux43eux439ux441ux442ux432ux43e-7.-ux43dux435ux43fux440ux435ux440ux44bux432ux43dux43eux441ux442ux44c-ux432ux435ux440ux43eux44fux442ux43dux43eux441ux442ux438-ux441ux43dux438ux437ux443}

Если

\[
A_1\subset A_2\subset\cdots,
\qquad
A=\bigcup_{n=1}^{\infty}A_n,
\]

то

\[
\mathbb P(A_n)\uparrow\mathbb P(A),
\qquad
\lim_{n\to\infty}\mathbb P(A_n)=\mathbb P(A).
\]

\subsubsection{Доказательство}\label{ux434ux43eux43aux430ux437ux430ux442ux435ux43bux44cux441ux442ux432ux43e-6}

Представим возрастающую последовательность через непересекающиеся слои:

\[
B_1=A_1,\qquad
B_n=A_n\setminus A_{n-1},\quad n\ge2.
\]

Тогда

\[
A_N=\bigsqcup_{n=1}^{N}B_n,
\qquad
A=\bigsqcup_{n=1}^{\infty}B_n.
\]

Следовательно,

\[
\mathbb P(A_N)=\sum_{n=1}^{N}\mathbb P(B_n),
\qquad
\mathbb P(A)=\sum_{n=1}^{\infty}\mathbb P(B_n).
\]

Переходя к пределу по \(N\), получаем

\[
\lim_{N\to\infty}\mathbb P(A_N)
=
\mathbb P(A).
\]

\(\square\)

\subsubsection{Свойство 8. Непрерывность вероятности
сверху}\label{ux441ux432ux43eux439ux441ux442ux432ux43e-8.-ux43dux435ux43fux440ux435ux440ux44bux432ux43dux43eux441ux442ux44c-ux432ux435ux440ux43eux44fux442ux43dux43eux441ux442ux438-ux441ux432ux435ux440ux445ux443}

Если

\[
A_1\supset A_2\supset\cdots,
\qquad
A=\bigcap_{n=1}^{\infty}A_n,
\]

то

\[
\mathbb P(A_n)\downarrow\mathbb P(A),
\qquad
\lim_{n\to\infty}\mathbb P(A_n)=\mathbb P(A).
\]

\subsubsection{Доказательство}\label{ux434ux43eux43aux430ux437ux430ux442ux435ux43bux44cux441ux442ux432ux43e-7}

Так как \(A_n\downarrow A\), дополнения возрастают:

\[
\overline{A_1}\subset \overline{A_2}\subset\cdots,
\qquad
\bigcup_{n=1}^{\infty}\overline{A_n}=\overline{A}.
\]

По непрерывности снизу:

\[
\mathbb P(\overline{A_n})\uparrow\mathbb P(\overline{A}).
\]

Используя формулу дополнения, получаем

\[
\mathbb P(A_n)=1-\mathbb P(\overline{A_n})
\downarrow
1-\mathbb P(\overline{A})
=
\mathbb P(A).
\]

\(\square\)

Для общих мер в этом утверждении нужно условие конечности хотя бы
\(\mu(A_1)<\infty\). Для вероятностной меры оно выполнено автоматически,
потому что \(\mathbb P(A_1)\le1\).

\subsubsection{Теорема 9. Функция распределения порождает
меру}\label{ux442ux435ux43eux440ux435ux43cux430-9.-ux444ux443ux43dux43aux446ux438ux44f-ux440ux430ux441ux43fux440ux435ux434ux435ux43bux435ux43dux438ux44f-ux43fux43eux440ux43eux436ux434ux430ux435ux442-ux43cux435ux440ux443}

Пусть \(F:\mathbb R\to[0,1]\) не убывает, непрерывна справа и

\[
\lim_{x\to-\infty}F(x)=0,\qquad
\lim_{x\to+\infty}F(x)=1.
\]

Тогда существует единственная вероятностная мера \(\mu_F\) на
\(\mathcal B(\mathbb R)\) такая, что

\[
\mu_F((-\infty,x])=F(x)
\]

для всех \(x\in\mathbb R\). Эквивалентно,

\[
\mu_F((a,b])=F(b)-F(a),\qquad a<b.
\]

\subsubsection{Комментарий к
доказательству}\label{ux43aux43eux43cux43cux435ux43dux442ux430ux440ux438ux439-ux43a-ux434ux43eux43aux430ux437ux430ux442ux435ux43bux44cux441ux442ux432ux443}

По программе здесь доказательство не требуется. Идея такая: сначала
задают массу полуинтервалов формулой

\[
\mu_F((a,b])=F(b)-F(a),
\]

затем задают массу конечных дизъюнктных объединений таких полуинтервалов
как сумму масс. Монотонность \(F\) гарантирует неотрицательность,
телескопические суммы дают конечную аддитивность, а непрерывность справа
нужна для счетной аддитивности на исходной алгебре. После этого
применяется теорема продолжения меры с алгебры или полуалгебры на
порожденную \(\sigma\)-алгебру. Поскольку полуинтервалы порождают
\(\mathcal B(\mathbb R)\), получается мера на борелевских множествах.
Единственность следует из того, что значения на порождающем классе
фиксированы.

\subsection{5. Примеры}\label{ux43fux440ux438ux43cux435ux440ux44b}

\subsubsection{Пример 1. Один бросок
монеты}\label{ux43fux440ux438ux43cux435ux440-1.-ux43eux434ux438ux43d-ux431ux440ux43eux441ux43eux43a-ux43cux43eux43dux435ux442ux44b}

Пусть

\[
\Omega=\{H,T\},\qquad
\mathcal F=2^\Omega.
\]

Для симметричной монеты:

\[
\mathbb P(\{H\})=\mathbb P(\{T\})=\frac12.
\]

Тогда

\[
\mathbb P(\Omega)=1,\qquad
\mathbb P(\varnothing)=0.
\]

Счетная аддитивность здесь фактически сводится к конечной, потому что
пространство конечно.

\subsubsection{Пример 2. Счетное пространство
исходов}\label{ux43fux440ux438ux43cux435ux440-2.-ux441ux447ux435ux442ux43dux43eux435-ux43fux440ux43eux441ux442ux440ux430ux43dux441ux442ux432ux43e-ux438ux441ux445ux43eux434ux43eux432}

Пусть

\[
\Omega=\mathbb N,\qquad
\mathcal F=2^{\mathbb N},
\]

и задана последовательность \(p_n\ge0\) такая, что

\[
\sum_{n=1}^{\infty}p_n=1.
\]

Тогда

\[
\mathbb P(A)=\sum_{n\in A}p_n
\]

задает вероятность на \(2^{\mathbb N}\). В частности,

\[
\mathbb P(\{n\})=p_n.
\]

Это типичная дискретная вероятностная мера.

\subsubsection{Пример 3. Вырожденное распределение и мера
Дирака}\label{ux43fux440ux438ux43cux435ux440-3.-ux432ux44bux440ux43eux436ux434ux435ux43dux43dux43eux435-ux440ux430ux441ux43fux440ux435ux434ux435ux43bux435ux43dux438ux435-ux438-ux43cux435ux440ux430-ux434ux438ux440ux430ux43aux430}

Пусть \(a\in\mathbb R\) и

\[
F(x)=
\begin{cases}
0, & x<a,\\
1, & x\ge a.
\end{cases}
\]

Тогда \(\mu_F=\delta_a\), то есть вся масса сосредоточена в точке \(a\):

\[
\mu_F(A)=
\begin{cases}
1, & a\in A,\\
0, & a\notin A.
\end{cases}
\]

Скачок функции распределения равен

\[
F(a)-F(a-)=1-0=1,
\]

поэтому

\[
\mu_F(\{a\})=1.
\]

\subsubsection{\texorpdfstring{Пример 4. Распределение Бернулли на
\(\mathbb R\)}{Пример 4. Распределение Бернулли на \textbackslash mathbb R}}\label{ux43fux440ux438ux43cux435ux440-4.-ux440ux430ux441ux43fux440ux435ux434ux435ux43bux435ux43dux438ux435-ux431ux435ux440ux43dux443ux43bux43bux438-ux43dux430-mathbb-r}

Пусть случайная величина принимает значения \(0\) и \(1\):

\[
\mathbb P(X=1)=p,\qquad
\mathbb P(X=0)=1-p.
\]

Ее функция распределения:

\[
F(x)=
\begin{cases}
0, & x<0,\\
1-p, & 0\le x<1,\\
1, & x\ge1.
\end{cases}
\]

У меры два атома:

\[
\mu_F(\{0\})=1-p,\qquad
\mu_F(\{1\})=p.
\]

\subsubsection{\texorpdfstring{Пример 5. Равномерное распределение на
\([0,1]\)}{Пример 5. Равномерное распределение на {[}0,1{]}}}\label{ux43fux440ux438ux43cux435ux440-5.-ux440ux430ux432ux43dux43eux43cux435ux440ux43dux43eux435-ux440ux430ux441ux43fux440ux435ux434ux435ux43bux435ux43dux438ux435-ux43dux430-01}

Функция распределения:

\[
F(x)=
\begin{cases}
0, & x<0,\\
x, & 0\le x<1,\\
1, & x\ge1.
\end{cases}
\]

Тогда для \(0\le a<b\le1\):

\[
\mu_F((a,b])=F(b)-F(a)=b-a.
\]

В этом случае мера совпадает с длиной на отрезке \([0,1]\).

\subsubsection{Пример 6. Экспоненциальное
распределение}\label{ux43fux440ux438ux43cux435ux440-6.-ux44dux43aux441ux43fux43eux43dux435ux43dux446ux438ux430ux43bux44cux43dux43eux435-ux440ux430ux441ux43fux440ux435ux434ux435ux43bux435ux43dux438ux435}

При параметре \(\lambda>0\):

\[
F(x)=
\begin{cases}
0, & x<0,\\
1-e^{-\lambda x}, & x\ge0.
\end{cases}
\]

Тогда для \(0\le a<b\):

\[
\mu_F((a,b])=e^{-\lambda a}-e^{-\lambda b}
=
\int_a^b \lambda e^{-\lambda t}\,dt.
\]

Плотность равна

\[
f(t)=\lambda e^{-\lambda t}\mathbf 1_{\{t\ge0\}}.
\]

\subsubsection{Пример 7. Почему конечной аддитивности
мало}\label{ux43fux440ux438ux43cux435ux440-7.-ux43fux43eux447ux435ux43cux443-ux43aux43eux43dux435ux447ux43dux43eux439-ux430ux434ux434ux438ux442ux438ux432ux43dux43eux441ux442ux438-ux43cux430ux43bux43e}

На алгебре конечных и ко-конечных подмножеств \(\mathbb N\) можно задать
функцию

\[
\nu(A)=
\begin{cases}
0, & A\text{ конечно},\\
1, & A\text{ ко-конечно}.
\end{cases}
\]

Она конечно-аддитивна и нормирована: \(\nu(\mathbb N)=1\). Но она не
является счетно-аддитивной, потому что

\[
\mathbb N=\bigsqcup_{n=1}^{\infty}\{n\},
\]

и тогда счетная аддитивность потребовала бы

\[
1=\nu(\mathbb N)=\sum_{n=1}^{\infty}\nu(\{n\})=0.
\]

Это показывает, что счетная аддитивность не является формальностью. Она
запрещает такие патологические вероятности и делает возможным предельный
анализ.

\subsection{6. Доказательства и
обоснования}\label{ux434ux43eux43aux430ux437ux430ux442ux435ux43bux44cux441ux442ux432ux430-ux438-ux43eux431ux43eux441ux43dux43eux432ux430ux43dux438ux44f}

\textbf{Почему \(\sigma\)-алгебра нужна вместе со счетной
аддитивностью.}

Счетная аддитивность говорит о вероятности счетного объединения:

\[
\mathbb P\left(\bigsqcup_{n=1}^{\infty}A_n\right).
\]

Чтобы левая часть имела смысл, само объединение должно быть событием.
Поэтому область определения вероятности должна быть замкнута
относительно счетных объединений. Именно это и обеспечивает
\(\sigma\)-алгебра. Если бы \(\mathcal F\) была только алгеброй, счетное
объединение событий могло бы не принадлежать \(\mathcal F\), и аксиома
счетной аддитивности не была бы применима.

\textbf{Почему непрерывность снизу эквивалентна счетной аддитивности при
конечной аддитивности.}

Если мера конечно-аддитивна и непрерывна снизу, то для попарно
непересекающихся \(A_n\) положим

\[
C_N=\bigsqcup_{n=1}^{N}A_n,
\qquad
C=\bigsqcup_{n=1}^{\infty}A_n.
\]

Тогда

\[
C_1\subset C_2\subset\cdots,
\qquad
C=\bigcup_{N=1}^{\infty}C_N.
\]

По непрерывности снизу:

\[
\mathbb P(C)=\lim_{N\to\infty}\mathbb P(C_N).
\]

По конечной аддитивности:

\[
\mathbb P(C_N)=\sum_{n=1}^{N}\mathbb P(A_n).
\]

Следовательно,

\[
\mathbb P(C)
=
\lim_{N\to\infty}\sum_{n=1}^{N}\mathbb P(A_n)
=
\sum_{n=1}^{\infty}\mathbb P(A_n).
\]

Значит получаем счетную аддитивность. Обратное направление уже доказано
в свойстве о непрерывности снизу. Поэтому в вероятностной мере счетная
аддитивность и непрерывность снизу тесно связаны.

\textbf{Почему правая непрерывность важна для функции распределения.}

Если \(\mu\) - вероятностная мера и

\[
F(x)=\mu((-\infty,x]),
\]

то \(F\) автоматически непрерывна справа. Действительно, если
\(x_n\downarrow x\), то множества

\[
(-\infty,x_n]\downarrow(-\infty,x].
\]

По непрерывности меры сверху:

\[
\mu((-\infty,x_n])\downarrow\mu((-\infty,x]).
\]

То есть

\[
F(x_n)\downarrow F(x).
\]

Это и есть правая непрерывность. Поэтому в обратной теореме о построении
меры по \(F\) правая непрерывность не случайное техническое условие, а
необходимое свойство любой настоящей функции распределения.

\textbf{Почему \(\mu_F(\{a\})=F(a)-F(a-)\).}

Рассмотрим интервалы

\[
\left(a-\frac1n,a\right],\qquad n=1,2,\ldots
\]

Они убывают к одноточечному множеству:

\[
\bigcap_{n=1}^{\infty}\left(a-\frac1n,a\right]=\{a\}.
\]

По непрерывности сверху:

\[
\mu_F(\{a\})
=
\lim_{n\to\infty}\mu_F\left(\left(a-\frac1n,a\right]\right).
\]

По определению меры на полуинтервалах:

\[
\mu_F\left(\left(a-\frac1n,a\right]\right)
=
F(a)-F\left(a-\frac1n\right).
\]

Переходя к пределу, получаем

\[
\mu_F(\{a\})=F(a)-F(a-).
\]

\textbf{Почему мера по функции распределения единственна.}

Полуинтервалы \((-\infty,x]\) порождают борелевскую \(\sigma\)-алгебру:

\[
\mathcal B(\mathbb R)=\sigma\{(-\infty,x]:x\in\mathbb R\}.
\]

Если две вероятностные меры совпадают на всех множествах
\((-\infty,x]\), то они совпадают на порожденной ими \(\sigma\)-алгебре.
Строго это доказывается через \(\pi\)-\(\lambda\)-теорему или через
теорему единственности продолжения меры. Поэтому функция распределения
однозначно задает борелевскую вероятность на прямой.

\subsection{7. Что написать на
доске}\label{ux447ux442ux43e-ux43dux430ux43fux438ux441ux430ux442ux44c-ux43dux430-ux434ux43eux441ux43aux435}

\begin{enumerate}
\def\labelenumi{\arabic{enumi}.}
\tightlist
\item
  Определение колмогоровской тройки:
\end{enumerate}

\[
(\Omega,\mathcal F,\mathbb P),\qquad
\mathbb P:\mathcal F\to[0,1],\qquad
\mathbb P(\Omega)=1.
\]

\begin{enumerate}
\def\labelenumi{\arabic{enumi}.}
\setcounter{enumi}{1}
\tightlist
\item
  Счетная аддитивность:
\end{enumerate}

\[
A_i\cap A_j=\varnothing,\ i\ne j
\quad\Rightarrow\quad
\mathbb P\left(\bigsqcup_{n=1}^{\infty}A_n\right)
=
\sum_{n=1}^{\infty}\mathbb P(A_n).
\]

\begin{enumerate}
\def\labelenumi{\arabic{enumi}.}
\setcounter{enumi}{2}
\tightlist
\item
  Базовые следствия:
\end{enumerate}

\[
\mathbb P(\varnothing)=0,\qquad
\mathbb P(\overline{A})=1-\mathbb P(A),
\]

\[
A\subset B\Rightarrow \mathbb P(A)\le\mathbb P(B),
\]

\[
\mathbb P(A\cup B)=\mathbb P(A)+\mathbb P(B)-\mathbb P(A\cap B),
\]

\[
\mathbb P\left(\bigcup_{n=1}^{\infty}A_n\right)
\le
\sum_{n=1}^{\infty}\mathbb P(A_n).
\]

\begin{enumerate}
\def\labelenumi{\arabic{enumi}.}
\setcounter{enumi}{3}
\tightlist
\item
  Непрерывность вероятности:
\end{enumerate}

\[
A_n\uparrow A\Rightarrow \mathbb P(A_n)\uparrow\mathbb P(A),
\]

\[
A_n\downarrow A\Rightarrow \mathbb P(A_n)\downarrow\mathbb P(A).
\]

\begin{enumerate}
\def\labelenumi{\arabic{enumi}.}
\setcounter{enumi}{4}
\tightlist
\item
  Функция распределения и порожденная мера:
\end{enumerate}

\[
F(x)=\mu_F((-\infty,x]),
\]

\[
\mu_F((a,b])=F(b)-F(a),
\]

\[
\mu_F(\{a\})=F(a)-F(a-).
\]

\begin{enumerate}
\def\labelenumi{\arabic{enumi}.}
\setcounter{enumi}{5}
\tightlist
\item
  Условия на \(F\):
\end{enumerate}

\[
F\text{ не убывает},\qquad
F\text{ непрерывна справа},
\]

\[
\lim_{x\to-\infty}F(x)=0,\qquad
\lim_{x\to+\infty}F(x)=1.
\]

\subsection{8. Типичные дополнительные
вопросы}\label{ux442ux438ux43fux438ux447ux43dux44bux435-ux434ux43eux43fux43eux43bux43dux438ux442ux435ux43bux44cux43dux44bux435-ux432ux43eux43fux440ux43eux441ux44b}

\textbf{Вопрос. Чем вероятностная мера отличается от обычной меры?}

Вероятностная мера - это мера полной массы \(1\). То есть она
счетно-аддитивна и неотрицательна, как обычная мера, но дополнительно
нормирована:

\[
\mathbb P(\Omega)=1.
\]

\textbf{Вопрос. Почему нельзя ограничиться конечной аддитивностью?}

Потому что в теории вероятностей постоянно встречаются предельные
события: счетные объединения, счетные пересечения, \(\limsup A_n\),
\(\liminf A_n\). Конечная аддитивность не дает непрерывности вероятности
по монотонным последовательностям событий. Без нее нельзя корректно
доказывать предельные теоремы.

\textbf{Вопрос. Счетная аддитивность применяется к любым \(A_n\)?}

Нет. Равенство с суммой применяется к попарно непересекающимся событиям.
Для произвольных событий верна только субаддитивность:

\[
\mathbb P\left(\bigcup_n A_n\right)\le \sum_n\mathbb P(A_n).
\]

\textbf{Вопрос. Почему \(\mathbb P(A)\le1\)?}

Потому что \(A\subset\Omega\), а по монотонности:

\[
\mathbb P(A)\le\mathbb P(\Omega)=1.
\]

\textbf{Вопрос. Что означает \(A_n\uparrow A\)?}

Это означает, что

\[
A_1\subset A_2\subset\cdots,
\qquad
A=\bigcup_{n=1}^{\infty}A_n.
\]

Тогда вероятности также возрастают к вероятности предела:

\[
\mathbb P(A_n)\uparrow\mathbb P(A).
\]

\textbf{Вопрос. Что означает \(A_n\downarrow A\)?}

Это означает, что

\[
A_1\supset A_2\supset\cdots,
\qquad
A=\bigcap_{n=1}^{\infty}A_n.
\]

Для вероятностной меры:

\[
\mathbb P(A_n)\downarrow\mathbb P(A).
\]

\textbf{Вопрос. Почему функция распределения должна быть непрерывна
справа?}

Потому что для любой меры

\[
F(x)=\mu((-\infty,x])
\]

и при \(x_n\downarrow x\) множества \((-\infty,x_n]\) убывают к
\((-\infty,x]\). Непрерывность меры сверху дает
\(F(x_n)\downarrow F(x)\).

\textbf{Вопрос. Что такое атом распределения?}

Атом в точке \(a\) - это положительная масса одноточечного множества:

\[
\mu(\{a\})>0.
\]

Для функции распределения:

\[
\mu(\{a\})=F(a)-F(a-).
\]

Значит атомы соответствуют скачкам функции распределения.

\textbf{Вопрос. Если функция распределения непрерывна, значит ли это,
что есть плотность?}

Нет. Непрерывность означает отсутствие атомов, но не гарантирует
существование плотности. Бывают непрерывные сингулярные распределения,
например распределение Кантора.

\textbf{Вопрос. Почему мера \(\mu_F\) единственна?}

Потому что значения на лучах \((-\infty,x]\) задают значения функции
распределения, а эти лучи порождают всю борелевскую \(\sigma\)-алгебру
\(\mathcal B(\mathbb R)\). Две вероятностные меры, совпадающие на таком
порождающем классе, совпадают на всех борелевских множествах.

\subsection{9. Типичные
ошибки}\label{ux442ux438ux43fux438ux447ux43dux44bux435-ux43eux448ux438ux431ux43aux438}

\begin{enumerate}
\def\labelenumi{\arabic{enumi}.}
\item
  Писать счетную аддитивность для непересекающихся событий не указав
  условие \(A_i\cap A_j=\varnothing\).
\item
  Путать аддитивность и субаддитивность. Для произвольных событий обычно
  верно только
\end{enumerate}

\[
\mathbb P\left(\bigcup_n A_n\right)\le\sum_n\mathbb P(A_n).
\]

\begin{enumerate}
\def\labelenumi{\arabic{enumi}.}
\setcounter{enumi}{2}
\item
  Называть любое подмножество \(\Omega\) событием. В абстрактной модели
  событие - это только элемент \(\mathcal F\).
\item
  Забывать нормировку \(\mathbb P(\Omega)=1\). Без нее была бы просто
  конечная мера, а не вероятностная мера.
\item
  Считать непрерывность сверху верной для любой меры без условий. Для
  общей меры нужно \(\mu(A_1)<\infty\). В вероятностном пространстве это
  условие автоматическое.
\item
  Писать для функции распределения меру интервала как \(F(b)-F(a)\) без
  уточнения типа интервала. Формула
\end{enumerate}

\[
\mu_F((a,b])=F(b)-F(a)
\]

соответствует правонепрерывной функции распределения
\(F(x)=\mu((-\infty,x])\).

\begin{enumerate}
\def\labelenumi{\arabic{enumi}.}
\setcounter{enumi}{6}
\item
  Думать, что непрерывная функция распределения обязательно имеет
  плотность. Это неверно: отсутствие скачков не равно абсолютной
  непрерывности.
\item
  Забывать, что теорема о мере, порожденной \(F\), строит меру на
  \(\mathcal B(\mathbb R)\), а не сразу на всех подмножествах
  \(\mathbb R\).
\end{enumerate}

\subsection{10. Связи с другими
билетами}\label{ux441ux432ux44fux437ux438-ux441-ux434ux440ux443ux433ux438ux43cux438-ux431ux438ux43bux435ux442ux430ux43cux438}

\textbf{Билет 01.} Здесь используются \(\sigma\)-алгебры и меры. Билет 2
добавляет счетную аддитивность и нормировку \(\mathbb P(\Omega)=1\).

\textbf{Билет 03.} Теорема о мере, порожденной функцией распределения,
доказывается через продолжение меры с полуалгебры или алгебры на
\(\sigma\)-алгебру.

\textbf{Билет 04.} Мера \(\mu_F\) строится на борелевской
\(\sigma\)-алгебре \(\mathcal B(\mathbb R)\), поэтому нужны борелевские
множества и пополнение меры.

\textbf{Билет 05.} Случайная величина - это измеримое отображение. Ее
распределение является образом вероятностной меры, а функция
распределения задает это распределение на прямой.

\textbf{Билет 06.} Распределения случайных величин и векторов опираются
на вероятностные меры и функции распределения.

\textbf{Билет 07.} Интеграл Лебега строится по мере. Математическое
ожидание - это интеграл по вероятностной мере.

\textbf{Билет 15 и Билет 16.} Теорема Колмогорова о согласованных
конечномерных распределениях и случайные процессы используют абстрактные
вероятностные пространства как базовую модель.

\subsection{11. Минимум на
удовлетворительно}\label{ux43cux438ux43dux438ux43cux443ux43c-ux43dux430-ux443ux434ux43eux432ux43bux435ux442ux432ux43eux440ux438ux442ux435ux43bux44cux43dux43e}

Нужно уметь сказать:

\begin{enumerate}
\def\labelenumi{\arabic{enumi}.}
\tightlist
\item
  Вероятностное пространство - это тройка
  \((\Omega,\mathcal F,\mathbb P)\).
\item
  \(\Omega\) - исходы, \(\mathcal F\) - \(\sigma\)-алгебра событий,
  \(\mathbb P\) - вероятность.
\item
  Главные условия:
\end{enumerate}

\[
\mathbb P(\Omega)=1,
\qquad
\mathbb P\left(\bigsqcup_{n=1}^{\infty}A_n\right)
=
\sum_{n=1}^{\infty}\mathbb P(A_n).
\]

\begin{enumerate}
\def\labelenumi{\arabic{enumi}.}
\setcounter{enumi}{3}
\tightlist
\item
  Из аксиом следуют \(\mathbb P(\varnothing)=0\),
  \(\mathbb P(\overline{A})=1-\mathbb P(A)\), монотонность.
\item
  Функция распределения \(F\) порождает меру \(\mu_F\) на
  \(\mathcal B(\mathbb R)\) по формуле
\end{enumerate}

\[
\mu_F((a,b])=F(b)-F(a).
\]

\begin{enumerate}
\def\labelenumi{\arabic{enumi}.}
\setcounter{enumi}{5}
\tightlist
\item
  Условия на \(F\): неубывание, правая непрерывность, пределы \(0\) и
  \(1\) на бесконечностях.
\end{enumerate}

\subsection{12. Минимум на
хорошо}\label{ux43cux438ux43dux438ux43cux443ux43c-ux43dux430-ux445ux43eux440ux43eux448ux43e}

Кроме минимума на удовлетворительно, нужно:

\begin{enumerate}
\def\labelenumi{\arabic{enumi}.}
\tightlist
\item
  Доказать \(\mathbb P(\varnothing)=0\) из счетной аддитивности.
\item
  Доказать формулу дополнения и монотонность.
\item
  Доказать счетную субаддитивность через дизъюнктизацию:
\end{enumerate}

\[
B_1=A_1,\qquad
B_n=A_n\setminus\bigcup_{k=1}^{n-1}A_k.
\]

\begin{enumerate}
\def\labelenumi{\arabic{enumi}.}
\setcounter{enumi}{3}
\tightlist
\item
  Сформулировать непрерывность вероятности снизу и сверху:
\end{enumerate}

\[
A_n\uparrow A\Rightarrow \mathbb P(A_n)\uparrow\mathbb P(A),
\]

\[
A_n\downarrow A\Rightarrow \mathbb P(A_n)\downarrow\mathbb P(A).
\]

\begin{enumerate}
\def\labelenumi{\arabic{enumi}.}
\setcounter{enumi}{4}
\tightlist
\item
  Объяснить, почему счетная аддитивность нужна для предельных событий.
\item
  Привести 2-3 примера мер, порожденных функцией распределения: Дирак,
  Бернулли, равномерное распределение.
\end{enumerate}

\subsection{13. Что добавить для
отлично}\label{ux447ux442ux43e-ux434ux43eux431ux430ux432ux438ux442ux44c-ux434ux43bux44f-ux43eux442ux43bux438ux447ux43dux43e}

Для сильного ответа стоит добавить:

\begin{enumerate}
\def\labelenumi{\arabic{enumi}.}
\tightlist
\item
  Связь счетной аддитивности с непрерывностью меры снизу: при конечной
  аддитивности непрерывность снизу фактически дает счетную аддитивность.
\item
  Контрпример конечно-аддитивной, но не счетно-аддитивной функции на
  конечных и ко-конечных подмножествах \(\mathbb N\).
\item
  Пояснение, что правая непрерывность функции распределения следует из
  непрерывности меры сверху.
\item
  Формулу атома:
\end{enumerate}

\[
\mu_F(\{a\})=F(a)-F(a-).
\]

\begin{enumerate}
\def\labelenumi{\arabic{enumi}.}
\setcounter{enumi}{4}
\tightlist
\item
  Объяснение, что непрерывная функция распределения не обязана иметь
  плотность.
\item
  Идею построения \(\mu_F\): полуинтервалы \(\to\) алгебра конечных
  объединений \(\to\) продолжение меры на \(\mathcal B(\mathbb R)\).
\item
  Упоминание, что доказательство существования \(\mu_F\) в этом билете
  не требуется, но относится к теореме продолжения меры.
\end{enumerate}

\subsection{14. Устная версия ответа на 5
минут}\label{ux443ux441ux442ux43dux430ux44f-ux432ux435ux440ux441ux438ux44f-ux43eux442ux432ux435ux442ux430-ux43dux430-5-ux43cux438ux43dux443ux442}

Вероятностное пространство в аксиоматике Колмогорова - это тройка
\((\Omega,\mathcal F,\mathbb P)\). Здесь \(\Omega\) - множество
элементарных исходов, \(\mathcal F\) - \(\sigma\)-алгебра событий, а
\(\mathbb P\) - вероятностная мера. Событиями являются не все
подмножества \(\Omega\), а только элементы \(\mathcal F\).

Вероятность удовлетворяет трем основным условиям: неотрицательность,
нормировка \(\mathbb P(\Omega)=1\) и счетная аддитивность. Счетная
аддитивность означает, что для попарно непересекающихся событий
\(A_1,A_2,\ldots\) выполнено

\[
\mathbb P\left(\bigsqcup_{n=1}^{\infty}A_n\right)
=
\sum_{n=1}^{\infty}\mathbb P(A_n).
\]

Это сильнее конечной аддитивности и нужно для работы с предельными
событиями. Например, события вида \(\bigcup_n A_n\), \(\bigcap_n A_n\) и
\(\limsup A_n\) естественно возникают в теории случайных процессов и
предельных теоремах.

Из счетной аддитивности выводятся основные свойства. Во-первых,
\(\mathbb P(\varnothing)=0\): применяем счетную аддитивность к
разложению \(\Omega=\Omega\sqcup\varnothing\sqcup\cdots\). Во-вторых, из
разложения \(\Omega=A\sqcup \overline{A}\) следует
\(\mathbb P(\overline{A})=1-\mathbb P(A)\). В-третьих, если
\(A\subset B\), то \(B=A\sqcup(B\setminus A)\), поэтому
\(\mathbb P(B)=\mathbb P(A)+\mathbb P(B\setminus A)\) и
\(\mathbb P(A)\le\mathbb P(B)\).

Также важна счетная субаддитивность:

\[
\mathbb P\left(\bigcup_n A_n\right)\le\sum_n\mathbb P(A_n).
\]

Она доказывается заменой \(A_n\) на непересекающиеся множества
\(B_1=A_1\), \(B_n=A_n\setminus\bigcup_{k<n}A_k\). Еще одно важное
следствие - непрерывность вероятности. Если \(A_n\uparrow A\), то
\(\mathbb P(A_n)\uparrow\mathbb P(A)\); если \(A_n\downarrow A\), то
\(\mathbb P(A_n)\downarrow\mathbb P(A)\).

Теперь пример меры, порожденной функцией распределения. Пусть
\(F:\mathbb R\to[0,1]\) не убывает, непрерывна справа и имеет пределы
\(0\) при \(x\to-\infty\) и \(1\) при \(x\to+\infty\). Тогда существует
единственная вероятностная мера \(\mu_F\) на \(\mathcal B(\mathbb R)\)
такая, что

\[
\mu_F((-\infty,x])=F(x).
\]

Отсюда для полуинтервала:

\[
\mu_F((a,b])=F(b)-F(a).
\]

Доказательство существования здесь обычно не требуется. Идея: задать
меру на полуинтервалах, затем на конечных объединениях полуинтервалов,
после чего применить теорему продолжения меры. Скачки \(F\) дают атомы:

\[
\mu_F(\{a\})=F(a)-F(a-).
\]

Например, ступенчатая функция распределения дает дискретную меру,
непрерывная линейная на \([0,1]\) функция дает равномерное
распределение, а функция \(1-e^{-\lambda x}\) при \(x\ge0\) дает
экспоненциальное распределение.

\subsection{15. Если осталось 2 минуты перед
экзаменом}\label{ux435ux441ux43bux438-ux43eux441ux442ux430ux43bux43eux441ux44c-2-ux43cux438ux43dux443ux442ux44b-ux43fux435ux440ux435ux434-ux44dux43aux437ux430ux43cux435ux43dux43eux43c}

Запомнить главное:

\[
(\Omega,\mathcal F,\mathbb P)
\]

\begin{itemize}
\tightlist
\item
  это пространство исходов, \(\sigma\)-алгебра событий и вероятность.
\end{itemize}

Главная аксиома:

\[
\mathbb P\left(\bigsqcup_{n=1}^{\infty}A_n\right)
=
\sum_{n=1}^{\infty}\mathbb P(A_n)
\]

для попарно непересекающихся \(A_n\).

Основные следствия:

\[
\mathbb P(\varnothing)=0,\qquad
\mathbb P(\overline{A})=1-\mathbb P(A),
\]

\[
A\subset B\Rightarrow\mathbb P(A)\le\mathbb P(B),
\]

\[
A_n\uparrow A\Rightarrow\mathbb P(A_n)\uparrow\mathbb P(A),
\]

\[
A_n\downarrow A\Rightarrow\mathbb P(A_n)\downarrow\mathbb P(A).
\]

Функция распределения \(F\) должна быть неубывающей, непрерывной справа,
с пределами \(0\) и \(1\) на бесконечностях. Она задает меру на
\(\mathcal B(\mathbb R)\):

\[
\mu_F((-\infty,x])=F(x),
\qquad
\mu_F((a,b])=F(b)-F(a).
\]

Скачок функции распределения равен массе точки:

\[
\mu_F(\{a\})=F(a)-F(a-).
\]

\newpage

\end{document}
